\documentclass[10pt,a4paper]{amsart}
\usepackage[margin=19mm]{geometry}
\usepackage{amsmath,amssymb,mathtools}
\usepackage[scr=rsfs,cal=euler]{mathalfa}
\usepackage{xfrac}
\usepackage[unicode,hidelinks,bookmarks=false]{hyperref}
\newcommand{\ie}{i.e.\ }
\newcommand{\cf}{cf.\ }
\newcommand{\resp}{respectively\ }
\newcommand{\Subsection}{\subsection}
\providecommand{\nobreakdash}{\nobreak\hbox{-}}
\setlength{\emergencystretch}{2em}
\multlinegap0pt
\usepackage{bm}
\usepackage{bbm}
\usepackage{dsfont}
\usepackage{xspace}
\usepackage{cancel}
\usepackage{mathtools}
\usepackage{amsthm}
\makeatletter
\@ifundefined{thm}{\newtheorem{thm}{Theorem}}{}
\@ifundefined{lem}{\newtheorem{lem}[thm]{Lemma}}{}
\makeatother
\newcommand\bC{{\mathbb C}}
\newcommand\bR{{\mathbb R}}
\newcommand\cB{{\mathcal B}}
\newcommand\cC{{\mathcal C}}
\newcommand\cL{{\mathcal L}}
\newcommand\cS{{\mathcal S}}
\newcommand\ve{\varepsilon}
\newcommand{\bbe}{\mathbbm e}
\newcommand{\bbm}{\mathbbm m}
\newcommand{\diam}{\operatorname{diam}}
\title[Central Limit Theorem for Sequential Dynamical Systems]{Central Limit Theorem for Sequential Dynamical Systems}
\author{Mark F. Demers, Carlangelo Liverani}
\date{Source: arXiv:2502.07765v2 (CC BY 4.0)}
\begin{document}
\maketitle

\noindent\emph{Source and licence.} Central Limit Theorem for Sequential Dynamical Systems. Version arXiv:2502.07765v2 (CC BY 4.0), distributed under the Creative Commons Attribution 4.0 International licence, as recorded in the arXiv licence metadata for this exact version and shown on the abstract page https://arxiv.org/abs/2502.07765v2. Full rights notice: licenses/arxiv-2502.07765-CC-BY-4.0.txt.\par\smallskip

\noindent\textit{Editorial scope.} Counted unit: the Thm whose source label is \texttt{thm:needed} in the article (source0.tex lines 1778--1792, source label \texttt{thm:needed}), delivered together with the article's own complete original proof of it (source0.tex lines 1793--1857, one \texttt{proof} environment of 65 lines). No mathematics is written, repaired or re-derived: statement and proof are the article's own text line for line. Required context retained uncounted: eq:archimedean (lines 1177--1180); eq:maybetoomuch (lines 1195--1197); lem:dual\_comp (lines 1285--1293); lem:upper\_dual (lines 1328--1332); eq:up\_lower (lines 1659--1661); eq:narrow bound (lines 1764--1768); eq:rad double (lines 1770--1773); eq:norms\_equiv (lines 4247--4249). Every definition, display and label of those lines is retained unchanged. Omitted from this package and not redistributed: 1--1176, 1181--1194, 1198--1284, 1294--1327, 1333--1658, 1662--1763, 1769--1769, 1774--1777, 1858--4246, 4250--4319. Nothing inside a retained excerpt is omitted or altered: every retained body is delivered exactly as the article's lines are written, so no omission receipt of any kind accompanies this package. Citations: the retained excerpts carry no citation command at all, so no bibliography entry of the article is transcribed and no reference is redistributed. Reference adaptation: none was needed and none was made; every reference inside the delivered unit resolves inside it, so no unresolved reference or citation is produced. Printed slips kept literal: no printed word, symbol, hypothesis or number of the counted statement or of its proof was altered. Layout and class: the article's own class is \texttt{amsart} with packages including amssymb, enumerate, xspace, esint, ulem, marvosym, bbm, todonotes, hyperref, geometry; the wrapper is the lane's pinned \texttt{amsart} editorial wrapper at 10pt on A4, which restates the theorem-like environments the retained text needs and adds the article's own macro definitions that the retained text needs (\texttt{\textbackslash bC}, \texttt{\textbackslash bR}, \texttt{\textbackslash bbe}, \texttt{\textbackslash bbm}, \texttt{\textbackslash cB}, \texttt{\textbackslash cC}, \texttt{\textbackslash cL}, \texttt{\textbackslash cS}, \texttt{\textbackslash diam}, \texttt{\textbackslash ve}), with extra wrapper packages (bm, bbm, dsfont, xspace, cancel, mathtools). The delivered numbering is the unit's own. Assets: no retained span contains a figure environment, a graphics-inclusion command, a PDF-page inclusion, a table or a TikZ picture; no image file is copied into this package, the package declares no asset dependency and no image file is redistributed with it. Licence: version arXiv:2502.07765v2 of this article is distributed under the Creative Commons Attribution 4.0 International licence, as recorded in the arXiv licence metadata for this exact version and shown on the abstract page https://arxiv.org/abs/2502.07765v2. A full rights notice is delivered at licenses/arxiv-2502.07765-CC-BY-4.0.txt. Characters: every retained excerpt is pure ASCII, so the delivered engine, XeLaTeX with its default Latin Modern text font, reports no missing character.
\begin{equation}
\label{eq:archimedean}
\mbox{for all $h\in  V$ there exists $\lambda\in\bR_+$ such that $\lambda\bbe-h\in C_\bR$.}
\end{equation}  

\begin{equation}\label{eq:maybetoomuch}
\bbm(h)\geq \kappa \|h\|.
\end{equation}

\begin{lem}\label{lem:dual_comp}
We have the following characterizations.
\begin{itemize}
  \item[a)] $\displaystyle \cC_\bC=\{h\in\cB_\bC\setminus\{0\}\;:\; \forall \ell, m\in \cC'_\bR,\; \Re\left(\ell(h)\overline{m(h)}\right)\geq 0\}$.
  \item[b)] $\displaystyle \cC_\bC'=\{ \ell\in \cB_{\bC}'\;:\; \forall x, y\in \cC_\bR,\; \Re(\ell(x)\overline{\ell(y)})>0 \}$. \smallskip
  \item[c)] $\displaystyle \cC'_\bC\supset\{\pm\ell\pm i p\;:\; \ell, p\in \mathring{\cC}'_\bR\}=:\hat\cC'_{\bC}$,
where $\mathring{\cC}'_\bR=\{\ell\in\cC'_{\bR}\;:\; \ell(x)>0 \, \forall x\in\cC_{\bR}\}$.
\end{itemize}
\end{lem}

\begin{lem}\label{lem:upper_dual} For each $\ell\in\cC_{\bC}'$ we have
\[
\|\ell \|' \leq \sqrt 2 |\ell(\bbe)|.
\]
\end{lem}

\begin{equation}\label{eq:up_lower}
\ell(\bbe_2)\|Ah\|e^{\Delta_{\bR}}\geq \ell(Ah)\geq \ell(\bbe_2)\kappa\|Ah\|e^{-\Delta_{\bR}}.
\end{equation}

\begin{equation}
\label{eq:narrow bound}
 \left|\frac{z \tilde\ell(A(x+iy)) \tilde p(\bbe_2)}{\tilde\ell(\bbe_2)z \tilde p(A(x+iy))}\right|\leq \sqrt 2\kappa^{-1}  c^{-1} e^{2\Delta_{\bR}} 
 \le 3\sqrt{2} \kappa^{-2} e^{4 \Delta_\bR} \, . 
\end{equation}

\begin{equation}
\label{eq:rad double}
\sup_{h\in\cC_{1, \bC}}\delta_{\cC_{2, \bC}}(A (h),\bbe_2)\leq \diam_{\delta_{\cC_{2,\bC}}}(A(\cC_{1,\bC}))\leq 2\sup_{h\in\cC_{1,\bC}}\delta_{\cC_{2, \bC}}(A h,\bbe_2),
\end{equation}

\begin{equation}\label{eq:norms_equiv}
\|h\|_{c,k}\leq \sqrt 2\|h\|_{r,k}.
\end{equation}

\begin{thm} \label{thm:needed} 
Let, for $i\in\{1,2\}$, $\cB_{i,\bR}$, $\cB_{i,\bC}$, $\cS_i$, $\bbe_i$ and $\bbm_i$ be as in the introduction to this section (in particular, the $\bbe_i$ and $\bbm_i$ satisfy equations \eqref{eq:archimedean} and \eqref{eq:maybetoomuch}, respectively).
Let $\cL\in L(\cB_{1,\bR},\cB_{2,\bR})$ and $\cL_\bC\in L(\cB_{1,\bC},\cB_{2,\bC})$.
Assume that $\cL(\cC_{1,\bR})\subset \cC_{2,\bR}$ and $\diam_H(\cL(\cC_{1,\bR})):=\Delta_\bR<\infty$. 
If there exists $\ve \in (0, \frac{\kappa^2}{12 \sqrt 2}e^{-2\Delta_\bR})$, such that
for all $\ell \in \cS_2$ and all $h \in \cC_{1,\bR}$,
\begin{equation}\label{eq:complex_bound}
\left|\ell(\cL_{\bC} h)-\ell(\cL h)\right|\leq \ve \ell(\cL h).
\end{equation}
Then $\cL_{\bC}(\cC_{1,\bC})\subset \cC_{2,\bC}$, and we have
\[
\diam_{\delta_{\cC_{2, \bC}}} (\cL_{\bC}(\cC_{1,\bC}))\leq 8\Delta_\bR +
 2 \ln[3\sqrt 2\kappa^{-2}]+ \tfrac{\sqrt{2}}{3} \kappa^2e^{-2\Delta_\bR} . 
\]
\end{thm}

\begin{proof}
For all $x,y\in \cC_{1,\bR}$ and $\ell\in\cS_2$ we have using \eqref{eq:complex_bound},
\begin{equation}\label{eq:real_comp_compa}
\begin{split}
\left|\ell(\cL_{\bC} (x+iy))-\ell(\cL (x+iy))\right|\leq& \left|\ell(\cL_{\bC} x)-\ell(\cL x)\right|
+\left|\ell(\cL_{\bC} y)-\ell(\cL y)\right|\\
\leq& \ve [\ell(\cL x)+ \ell(\cL y)]\leq \ve  \sqrt 2 |\ell(\cL (x+iy))|.
\end{split}
\end{equation}
It follows that for all $h\in\cC_{1,\bC}$ and $\ell,p\in\mathring\cC_{2,\bR}'$, provided $\ve\leq \frac{1}{\sqrt 2}$,
\[
\begin{split}
\Re(\ell(\cL_{\bC} h)\overline{p(\cL_{\bC} h)})
& =
\Re\Bigg(\ell(\cL h)\overline{p(\cL h)}  + \ell((\cL_{\bC}-\cL) h )\overline{p((\cL_{\bC}-\cL) h )}  \\
& +\ell((\cL_{\bC}-\cL) h)\overline{p(\cL h )}+\ell(\cL h )\overline{p((\cL_{\bC}-\cL) h )} \Bigg)\\
&\geq \Re\left(\ell(\cL h )\overline{p(\cL h)}\right)
- 3 \sqrt 2\ve|\ell(\cL h )||{p(\cL h )}|\\
&\geq \Re\left(\ell(\cL h )\overline{p(\cL h)}\right)
-6\sqrt 2\ve \ell(\bbe_2)p(\bbe_2)\|\cL h\|\|\cL h \|
\end{split}
\]
where, in the first inequality we have used \eqref{eq:real_comp_compa} and in the second inequality we have
used Lemma \ref{lem:dual_comp}(c) and Lemma~\ref{lem:upper_dual}. 
Since we can assume $h=x+iy$ with $x,y\in \cC_{1,\bR}$, recalling \eqref{eq:up_lower} we have
\[
\begin{split}
\Re\left(\ell(\cL h )\overline{p(\cL h )}\right) & = \ell(\cL(x))p(\cL x )+\ell(\cL y )p(\cL y)\\
& \geq \ell(\bbe_2)p(\bbe_2)\kappa^2 e^{-2\Delta_\bR} \left( \|\cL x\|^2 + \|\cL y \|^2 \right) \\
& \geq  2^{-1} \ell(\bbe_2)p(\bbe_2)\kappa^2 e^{-2\Delta_\bR} \| \cL h \|^2 \, , 
\end{split}
\]
where we have used \eqref{eq:norms_equiv}.  Hence, $\Re(\ell(\cL_{\bC}(h))\overline{p(\cL_{\bC}(h))})>0$, provided $\ve < \frac{\kappa^2}{ 12 \sqrt 2}e^{-2\Delta_\bR}$. 

Note that if $\ell \in \cC_{2,\bR}'$, then $\ell+\alpha\bbm\in\mathring\cC_{2,\bR}'$ for each $\alpha>0$.  Thus the above implies $\Re(\ell(\cL_{\bC}(h))\overline{p(\cL_{\bC}(h))}) \ge 0$ for
all $\ell, p \in \cC_{2, \bR}'$. 
Applying Lemma~\ref{lem:dual_comp}(a), 
 we have $\cL_{\bC}(\cC_{1,\bC})\subset \cC_{2,\bC}$.

Finally, for all $\ell,p\in\cC_{2,\bR}'$ and $h\in\cC_{1,\bC}$, it follows by \eqref{eq:real_comp_compa} that,
\[
\begin{split}
\left|\frac{\ell(\cL_{\bC}h)p(\cL h)}{\ell(\cL h)p( \cL_{\bC}h)}\right|&\leq 
\frac{|\ell(\cL h)p(\cL h)|+|\ell([\cL_{\bC}- \cL  ]h)p(\cL h)|}{|\ell(\cL h)p( \cL h)|
-|\ell(\cL h)p([\cL_{\bC}- \cL  ]h)|}\\
&\leq \frac{1+  \sqrt 2\ve}{1- \sqrt 2\ve}  \leq 1 + 4 \ve, 
\end{split}
\]
since we assumed $\ve < \frac 1{ 12 \sqrt{2}}$.


The above implies 
\begin{equation}
\label{eq:needed later}
\begin{split}
\delta_{\cC_{2, \bC}}(\cL_{\bC}h,\bbe_2)&=\sup_{\ell,p}\ln \frac{|\ell(\cL_{\bC}h)p(\bbe_2)|}{|\ell(\bbe_2)p( \cL_{\bC}h)|} \\
&\leq \sup_{\ell,p}\ln \frac{|\ell(\cL h)p(\bbe_2)|}{|\ell(\bbe_2)p( \cL h)|}+\sup_{\ell,p}\ln \frac{|\ell(\cL_{\bC}h)p(\cL h)|}{|\ell(\cL h)p( \cL_{\bC}h)|}\\
&\leq \delta_{\cC_{2, \bC} }(\cL h,\bbe_2)+ 4 \ve
\leq 4\Delta_\bR +  \ln[3 \sqrt 2\kappa^{-2}] + 4  \ve \\
&\leq 4\Delta_\bR +   \ln[3\sqrt 2\kappa^{-2}]+   \tfrac{1}{ 3 \sqrt{2} }  \kappa^2e^{-2\Delta_\bR},
\end{split}
\end{equation}
where, in the next to last inequality, we have applied \eqref{eq:narrow bound}. We conclude using
\eqref{eq:rad double}.
\end{proof}

\end{document}
